\scanpage{630}
\mspage{18}
adjoints $(u^*,u_*)$, cette relation équivaut à

\[
TC'\unsure{\longleftrightarrow|\Longleftrightarrow}u_*(F).
\]

Comme $TC'\overset{\mathrm{def}}{=}u_!(TC)^{\sim}$, ceci équivaut à

\[
u_!(TC)\longleftrightarrow u_*(\unsure{F}),
\]

ou encore (en réappliquant 1.5 à $u_!(TC),F$
au lieu de $(TC',F)$), à

\[
u^*u_!(TC)\longleftrightarrow F\qquad\text{i.e.}\quad
u^*u_!(TC)\subset\underset{\substack{\parallel\\ F^*}}{TC},
\]

or c'est là l'hypothèse \deleted{\unsure{a)}} b'').

\uline{Cas B} Il suffit encore de prouver
2.1.7. Or, du fait que $u_!(TC)$ \deleted{\unsure{admet}}\inserted{\unsure{admet}} \unsure{est}
\unsure{un} \tikz[baseline=(under.base)]{\node[inner sep=0pt] (under) {\deleted{\unsure{petit ensemble} $Q$-générateur}}; \draw[dashed] ([xshift=.6em,yshift=-.2ex]under.south west) -- ([xshift=-.3em,yshift=-.2ex]under.south east);} à gauche,
$u_!(TC)^{\sim}$ est la clôture de $u_!(TC)$ pour les
opérations de cochgt. de base, de composition
transfinie, de facteur direct. Comme
$u^*$ commute à \unsure{ttes|toutes} ces opérations (comme
il commute aux \inserted{petites} $\unsure{\underline{\lim}|\varinjlim}$), et \deleted{\unsure{étant donn}} comme
$TC$ est dans $\mathrm{Fl}\,M$ est stable par ces opérations,
on voit donc que pour montrer qu'\unsure{on|en}
$u^{\unsure{*|'}}$ applique $u_!(TC)^{\sim}$ dans $TC$, il suffit de
voir qu'il applique $u_!(TC)$ dans $TC$, ce qui
est encore l'hypothèse b'').

\uline{Cas C} On va vérifier b) sous la forme
équivalente b'). Soit donc $f'\in\mathrm{Fl}\,M'$,
\deleted{\unsure{factorisons} \unsure{u} \unsure{soit}} $f':X'\unsure{\text{---}|\longrightarrow}Y^{\unsure{F|'|*}}$,
% archive sheet number: 629
